In this section, we introduce the statistical metrics used to characterise the precision of our photometric redshifts. We then assess the accuracy of the point estimates and evaluate the performance of the conditional‐density modelling, thereby validating our data‐augmentation technique via comparative analysis.

\subsection{Statistical metrics}
We quantify the systematic error of each photometric‐redshift estimate \(z_\mathrm{phot}\) relative to its true redshift \(z_\mathrm{true}\) by defining the scaled bias
\begin{equation} \label{eqa:4}
    \delta_z = \frac{z_\mathrm{phot} - z_\mathrm{true}}{1 + z_\mathrm{true}}.
\end{equation}
From the distribution of \(\delta_z\), we derive two robust metrics for any galaxy subset. First, the median scaled bias, \(\bar{\delta}_z \equiv \mathrm{median} \left( \delta_z \right)\), is adopted in preference over the mean to mitigate sensitivity to outliers. Second, the Normalised Median Absolute Deviation (NMAD) of \(\delta_z\), denoted \(\sigma_z\), is defined as
\begin{equation} \label{eqa:5}
    \sigma_z = 1.4826 \times \mathrm{median} \left( \left| \delta_z - \bar{\delta}_z \right| \right).
\end{equation} 
The factor \(1.4826\) rescales the median absolute deviation into an estimator equivalent to the standard deviation for a Gaussian distribution.

Furthermore, we consider the proportion at which galaxies fall outside the outlier boundary \(\left| \delta_z \right| = 0.15\), as defined in previous studies \citep[e.g.][]{2012ApJS..198....1F, 2014ApJS..215...27Y, 2017A&A...605A..70D, 2017A&A...608A...3B, 2019A&A...622A...3D}, termed outlier fraction, 
\begin{equation} \label{eqa:6}
    f_{\mathrm{o}} = \frac{N_\mathrm{g} \left( \left| \delta_z \right| > 0.15 \right)}{N_\mathrm{g}},
\end{equation} 
where \(N_{\mathrm{g}}\) represents the total count of a specific galaxy subset. In addition, we incorporate the rates of catastrophic failures where \(\left| z_{\mathrm{phot}} - z_{\mathrm{true}} \right| > 1.0\), as outlined in \cite{2024ApJ...967L...6M}, denoted as
\begin{equation} \label{eqa:7}
    r_{\mathrm{c}} = \frac{N_\mathrm{g} \left( \left| z_{\mathrm{phot}} - z_{\mathrm{true}} \right| > 1.0 \right)}{N_\mathrm{g}}.
\end{equation}
These metrics evaluate the prevalence of notable systematic errors in the photometric-redshift determinations. In addition to point estimates, we evaluate the performance of the photometric‐redshift conditional density modelling, \(p \left(z | \mathcal{S}_n, \mathcal{P}_n \right) \). For simplicity, the explicit dependence on the \(n\)th spectroscopic dataset \(\mathcal{S}_n\) and photometric dataset \(\mathcal{P}_n\) will henceforth be omitted. We employ the cumulative distribution function,
\begin{equation} \label{eqa:8}
    q \left(z\right) \equiv \int_{0}^{z} p \left(z^\prime\right) \,\mathrm{d}z^\prime,
\end{equation}
to define summary statistics owing to its well‐behaved, bounded range. The Probability‐Integral-Transform (PIT) value -- defined as the CDF evaluated at the true redshift \(z_{\mathrm{true}}\), \(q \left( z_{\mathrm{true}} \right)\) -- has been widely used in the literature \citep[e.g.][]{2012MNRAS.421.1671B, 2017MNRAS.468.4556F, 2018PASJ...70S...9T, 2020MNRAS.499.1587S}. For an ideal conditional density model, the PIT values of the full sample follow a uniform distribution. We use PIT histograms to compare conditional-density modelling across the selected galaxy samples.

The normalised PIT histograms, \( \mathcal{H}_k \left[ q \left( z_\mathrm{true} \right) \right] \) for \(k = 1, \ldots, K\), are calculated across \(K = 10\) equally sized quantile bins, ranging from \(\left[0.0, 1.0\right]\), with intervals of \(0.1\). The ordinate \(\phi[q(z_\mathrm{true})]\) in Figure~\ref{fig:10} denotes this histogram density, whose height in bin \(k\) is \(\mathcal{H}_k\). We calculate these histograms for both the full lens and source samples as well as for each redshift interval defined below. Thus, for each normalised PIT histogram, we then define the divergence loss, \(\mathcal{D}_q\), as the Root-Mean-Square deviation of the normalised PIT histogram from uniformity:
\begin{equation} \label{eqa:9}
    \mathcal{D}_q = \sqrt{\frac{1}{K} \sum_{k = 1}^{K} \left\{ \mathcal{H}_k \left[ q \left( z_\mathrm{true} \right) \right] - 1 \right\}^2},
\end{equation}
where unity denotes the value of the normalised histogram for an ideal uniform distribution. We use this metric as an empirical indicator of conditional-density modelling performance across the selected galaxy samples.

For the lens sample, each metric is assessed across five redshift intervals, each with a width of \(0.3\), covering the range \(0.0 < z_\mathrm{true} < 1.5\). Similarly, for the source sample, six intervals of width \(0.5\) are used, encompassing \(0.0 < z_\mathrm{true} < 3.0\). The outcomes of our statistical metrics are depicted as a function of true redshift in Figure~\ref{fig:8}, with comprehensive discussions provided thereafter.

It is important to note that this interval binning is solely used for assessing the performance of the photometric redshifts, rather than for tomographic binning in cosmological analysis. The latter is based on the derived point estimates and will be addressed in our subsequent paper.

\begin{figure*}
    \centering
    \begin{subfigure}[b]{0.48\linewidth}
        \centering
        \includegraphics[width=\linewidth]{Figures/Y1/METRIC.pdf}
        \caption{LSST Y1}
    \end{subfigure}
    \hspace{0.0\linewidth}
    \begin{subfigure}[b]{0.48\linewidth}
        \centering
        \includegraphics[width=\linewidth]{Figures/Y10/METRIC.pdf}
        \caption{LSST Y10}
    \end{subfigure}
    \caption{Statistical metrics plotted as a function of true redshift, with LSST Y1 in the left panel and LSST Y10 in the right. Within each panel, the metrics are ordered from top to bottom: the median scaled bias \(\bar{\delta}_z\); the normalised median absolute deviation (NMAD) of the scaled bias \(\sigma_z\); the outlier fraction \(f_{\mathrm{o}}\); the catastrophic failure rate \(r_{\mathrm{c}}\); and the divergence loss of the normalised PIT histogram \(\mathcal{D}_q\). In each row, the left and right subplots correspond to the lens and source samples, respectively. Filled markers denote the median value of each metric across all experiments, with error bars indicating the central \(68\%\) interval across experiments. Green curves show in-sample results from models trained and evaluated on the same \texttt{Combination} datasets; orange curves correspond to \texttt{Application} datasets trained on these augmented \texttt{Combination} datasets; and blue curves represent \texttt{Application} datasets trained on the unaugmented \texttt{Degradation} datasets. The grey bands show SRD-motivated comparison benchmarks of \(0.02\) for lenses and \(0.03\) for sources; these are not ensemble-redshift calibration requirements.}
    \label{fig:8}
\end{figure*}

\subsection{Evaluation of point estimate accuracy}
Figure~\ref{fig:9} provides a direct visual assessment of the performance of our photometric redshift model under the baseline experiment. As depicted in the top panels of Figure~\ref{fig:9}, the in-sample evaluation on the \texttt{Combination} dataset shows excellent agreement between \(z_{\mathrm{phot}}\) and \(z_{\mathrm{true}}\), indicating that our photometric redshift model reliably captures the colour–redshift relation. By selecting the peak of the conditional PDF, we achieve negligible bias for both the bright lens sample (left panels) and the fainter source sample (right panels), even at \(z_\mathrm{true} \gtrsim 1.5\). The resulting mean absolute values of the scaled biases, \(\overline{\left|\delta_z\right|}\), lie below the SRD-motivated comparison benchmarks over almost the entire redshift range. These results demonstrate the accuracy of \texttt{FlexZBoost} on the fitted \texttt{Combination} sample.

Upon utilising the complete photometric \texttt{Application} datasets, the central panels of Figure~\ref{fig:9} demonstrate that the lens galaxies maintain a similar level of performance. In the bottom-left subplots for both the Y1 and Y10 \texttt{Application} lens samples, the mean absolute values of the scaled biases, \(\overline{\left|\delta_z\right|}\), highlighted by the black dashed lines, are notably beneath the SRD-motivated comparison benchmarks, depicted by black dotted lines. This signifies an exceptionally accurate model for this luminous population. Furthermore, practically no outliers or catastrophic failures occur in either instance, reinforcing the robustness of our photometric-redshift point estimates for bright galaxies and thus supporting precision cosmology through large-scale-structure analyses. In contrast, for the source samples, \(\overline{\left|\delta_z\right|}\) remains close to the comparison benchmark up to \(z_\mathrm{true} \lesssim 1.5\). Nonetheless, a considerable residual bias becomes evident at \(z_\mathrm{true} \gtrsim 1.5\).

Similar trends are evident in Figure~\ref{fig:8}, where the green and orange curves both correspond to models trained on the \texttt{Combination} datasets but evaluated on the \texttt{Combination} and \texttt{Application} datasets, respectively. This comparison shows that, for the lens sample, the median scaled bias \(\bar{\delta}_z\) and its NMAD \(\sigma_z\) generally remain within the grey comparison bands, with the exception of the bias in the highest true-redshift Y1 lens interval. For the source sample, however, we observe a significant increase in \(\bar{\delta}_z\), accompanied by a rise in its NMAD, as well as marked increases in the outlier fraction \(f_{\mathrm{o}}\), and the catastrophic failure rate \(r_{\mathrm{c}}\).

This suggests that the machine‐learning model struggles to extrapolate beyond its training domain in colour–SED space. Contributing factors include photometric noise -- which perturbs the mapping between observed SEDs and redshift, particularly for faint galaxies, thereby introducing uncertainty into \(z_{\rm phot}\) -- and systematic mismatches in SED properties. Specifically, the redshift upper limit of the \texttt{Degradation} dataset compels the model to rely on subtly discrepant galaxies in the augmented training dataset -- enhanced with high-redshift \texttt{CosmoDC2} galaxies -- compared to those in the photometric \texttt{Application} dataset from \texttt{OpenUniverse2024}. As discussed in Section~\ref{sec:3}, this SED discrepancy constitutes a primary source of bias in our photometric-redshift modelling.

The bottom two panels of Figure~\ref{fig:9} provide a qualitative comparison of photometric–redshift estimates trained exclusively on the \texttt{Degradation} datasets without augmentation. For the lens samples, incorporating the \texttt{Augmentation} datasets yields only marginal improvements: the mean absolute scaled biases remain effectively unchanged whether or not augmentation is applied. Nonetheless, for the source sample in the Y1 scenario, training on the unaugmented \texttt{Degradation} dataset results in virtually no galaxies with \(z_{\mathrm{phot}} \gtrsim 1.5\). A quantitative comparison of the orange and blue curves in Figure~\ref{fig:8} reinforces these conclusions. For the lens sample, these curves coincide nearly perfectly across all statistical metrics in both the LSST Y1 and Y10 cases. This reflects that the brightest galaxies are already well represented in the spectroscopic training sets.

Conversely, not employing data augmentation for the source sample results in significantly high outlier fractions (\(f_{\mathrm{o}} \gtrsim 0.5\)) and catastrophic failure rates (\(r_{\mathrm{c}} \gtrsim 0.5\)) at \(z_{\mathrm{true}} \gtrsim 2.0\) in the LSST Y1 scenario. Introducing SOM‐based data augmentation reduces the outlier fraction by approximately \(10 \,{-}\, 20\%\), lowers the catastrophic failure rate by a factor of \(\sim 2\), and decreases the median scaled bias by a factor of \(\sim 2 \,{-}\, 3\), thereby demonstrating the efficacy of augmentation in mitigating both bias and outliers in LSST Y1 photometric redshift point estimates. Although the NMAD of the scaled bias \(\sigma_z\) increases after data augmentation, this does not indicate a degradation in photometric performance. On the contrary, the low NMAD observed without augmentation is driven by a high fraction of catastrophic failures, which cluster galaxies around large bias values as seen in Figure~\ref{fig:9}. Consequently, the rise in NMAD post‐augmentation reflects the suppression of these extreme biases. Furthermore, the pronounced strip-like features evident in Figure~\ref{fig:9} after augmentation reflect intrinsic degeneracies in galaxy SEDs: in the absence of representative high-redshift training samples, such galaxies would otherwise be classified as catastrophic failures and, consequently, vanish from the regime around one-to-one sequence.

In the LSST Y10 scenario, the \texttt{Degradation} datasets comprise deeper spectroscopic observations than in Y1, thereby yielding improved photometric redshift performance for source galaxies even without augmentation. Nonetheless, training with the \texttt{Augmentation} datasets still confers benefits: both the median scaled bias \(\bar{\delta}_z\) and the catastrophic failure rate \(r_{\mathrm{c}}\) decrease by up to a factor of two in the highest-redshift interval (\(2.5 < z_{\mathrm{true}} < 3.0\)). However, in the Y10 case the gains are not uniform across all metrics. For instance, between \(1.5 < z_{\mathrm{true}} < 2.5\), the outlier fraction \(f_\mathrm{o}\) increases, accompanied by modest rises in the median scaled bias and its NMAD. This is likely due to residual mismatches in galaxy SEDs, reflecting limitations of our current augmentation implementation. As shown in Figure~\ref{fig:1}, the divergence between the \texttt{CosmoDC2} and \texttt{OpenUniverse2024} SEDs becomes more pronounced for the fainter galaxies probed at Y10 depths, which can introduce ambiguity into photometric-redshift training and degrade performance. This underscores the necessity for enhancing our augmented mock catalogues in preparation for future LSST Y10 photometric redshift analyses. 

Moreover, for both Y1 and Y10 scenarios, the large error bars on the blue curves denote substantial scatter in performance across unaugmented experiments -- highlighting sensitivity to the spectroscopic selection function -- whereas these variations are significantly reduced in the orange curves once augmentation is applied. Although the augmentation is not flawless, this reduced scatter helps demonstrate the efficacy of data augmentation.

\begin{figure*}
    \centering
    \begin{subfigure}[b]{0.48\linewidth}
        \centering
        \includegraphics[width=\linewidth]{Figures/Y1/POINT.pdf}
        \caption{LSST Y1}
    \end{subfigure}
    \hspace{0.0\linewidth}
    \begin{subfigure}[b]{0.48\linewidth}
        \centering
        \includegraphics[width=\linewidth]{Figures/Y10/POINT.pdf}
        \caption{LSST Y10}
    \end{subfigure}
    \caption{Direct visual comparison of \(z_{\mathrm{phot}}\) versus \(z_{\mathrm{true}}\) for the baseline experiment: LSST Y1 (left) and Y10 (right). From top to bottom, each row corresponds to (i) the \texttt{Combination} dataset used as an in-sample reference; (ii) the \texttt{Application} dataset trained on the augmented \texttt{Combination} dataset; and (iii) the \texttt{Application} dataset trained on the \texttt{Degradation} dataset without data augmentation. For every row, the upper panels show two-dimensional histograms for the lens (left) and source (right) samples, where the colour scale represents galaxy counts. The solid black line marks the ideal one‐to‐one relation \(z_{\mathrm{phot}} = z_{\mathrm{true}}\), and the dot‐dashed lines indicate the outlier threshold \(\left|\delta_z\right| = 0.15\). The lower panels show the absolute scaled bias \(\left|\delta_z\right|\) as a function of \(z_{\mathrm{true}}\) for lens (left) and source (right) galaxies. The dashed black line denotes the mean absolute scaled bias \(\overline{\left|\delta_z\right|}\), while the dotted lines show comparison benchmarks of \(0.03\) for lenses and \(0.05\) for sources, motivated by the LSST DESC SRD scatter assumptions.}
    \label{fig:9}
\end{figure*}

\subsection{Assessment of conditional density modelling}
In parallel with our point‐estimate analysis, we also evaluate the quality of conditional density modelling. In the top panels of Figure~\ref{fig:10}, the normalised PIT histograms for the \texttt{Combination} datasets demonstrate the in-sample conditional‐density performance. For the lens samples, in both the Y1 and Y10 scenarios, a central peak exceeding unity is observed, along with dips falling below the flat line near \(q \left( z_\mathrm{true} \right) = 0\) and \(1\), across all subsets at varying redshifts. For the full lens samples, this suggests that the conditional-PDF widths for these bright populations are slightly overestimated, resulting in an excess of objects at intermediate quantiles. Such over‐smoothing may stem from the basis‐function decomposition of the density model, since photometric redshift estimates for the low-redshift brightest galaxies would otherwise be expected to be sharply constrained.

The normalised PIT histogram of all source galaxies, as shown in the solid black curve, reveals elevated values at the quantile boundaries, whilst the central peak approximates unity. Comparable patterns can be observed in histograms of low-redshift sources (\(z_\mathrm{true} \lesssim 1.5\)), which suggests their predominance in the overall distribution. For the full source sample, this pattern suggests that some conditional PDFs are too narrow or systematically biased, unlike the bright lens samples. In the high-redshift region (\(z_\mathrm{true} \gtrsim 1.5\)), the histogram presents a marked slope even for this in-sample reference, alongside the point-estimate biases shown in Figures~\ref{fig:8} and \ref{fig:9}. These trends are also demonstrated in the increasing divergence loss of normalised PIT histograms in the bottom row of Figure~\ref{fig:8}, which correspond to the rise in point-estimate outliers and catastrophic failures, highlighting the difficulty of modelling precise PDFs at high redshifts.

In the central panels of Figure~\ref{fig:10}, consistent with our point‐estimate results, applying the photometric‐redshift model to the \texttt{Application} datasets leaves the lens‐sample PIT distributions largely unchanged but precipitates a pronounced deterioration in the source‐sample performance. In particular, at \(z\gtrsim1.5\), the slope of the PIT histogram becomes markedly steeper, consistent with the increased point-estimate bias for faint galaxies. This trend is also reflected in the upward shift of the divergence‐loss curves from green (\texttt{Combination} datasets) to orange (\texttt{Application} datasets) in Figure~\ref{fig:8}. These findings emphasise that the inclusion of faint, high-redshift galaxies -- coupled with photometric noise and non-representative SEDs in the training data -- exacerbates the challenges of conditional-density modelling, especially for high redshift galaxies. This is one of the primary challenges in calibrating ensemble redshift distributions and will be addressed in detail in our forthcoming paper, where we will explore how various estimators can be combined to mitigate systematic biases at the population-level.

Comparison with models trained solely on the unaugmented \texttt{Degradation} dataset further emphasises the benefits of data augmentation for conditional‐density modelling. For the lens sample, augmentation induces a negligible change. In contrast, for the source sample under the LSST Y1 scenario, the divergence loss of the normalised PIT histogram decreases by approximately \(25 \,{-}\, 50\%\) even at low redshift (\(z_{\rm true} \lesssim 1.5\)), whereas improvements in point‐estimate metrics remain largely confined to high‐redshift galaxies. This effect is evident in the lower panels of Figure~\ref{fig:10}. With augmented training, the edge peaks of the normalised PIT histograms in the low‐redshift bins are moderately suppressed and the central values approach unity, yielding distributions that are closer to uniform. This behaviour is consistent with a more faithful representation of SED degeneracies, owing to the inclusion of high‐redshift galaxies in the photometric‐redshift training.

At high redshift (\(z_{\rm true}\gtrsim1.5\)) in the LSST Y1 scenario, the unaugmented PIT histograms collapse to near zero except for a spike at \(q(z_{\rm true})=1\), indicative of pronounced systematic biases, consistent with the trends in Figures~\ref{fig:8} and \ref{fig:9}. Training on the augmented \texttt{Combination} datasets still produces a tilt relative to a uniform distribution, but a substantial fraction of galaxies exhibit markedly reduced bias, resulting in a non‐catastrophic normalised PIT histogram. This observation aligns with the reduction in catastrophic failure rates shown in Figure~\ref{fig:8}, highlighting that augmentation effectively restores missing information in the training dataset.

In the LSST Y10 scenario, the reduction in divergence loss for low-redshift galaxies remains pronounced, by up to a factor of \(\sim 2\) around \(z_\mathrm{true} \sim 1.5\), providing empirical support for improved conditional-density modelling. At higher redshifts, however, only modest gains are seen. This is partly because the deeper spectroscopic \texttt{Degradation} datasets expected for LSST Y10 analysis already enhance the conditional-density estimates for high-redshift galaxies even without augmentation. Additionally, increasingly severe SED mismatches between the \texttt{Augmentation} and \texttt{Application} datasets for faint, high-redshift galaxies limit further improvements. Nevertheless, data augmentation still reduces the scatter in photometric-redshift performance arising from variations in the spectroscopic selection function.

\begin{figure*}
    \centering
    \begin{subfigure}[b]{0.48\linewidth}
        \centering
        \includegraphics[width=\linewidth]{Figures/Y1/DENSITY.pdf}
        \caption{LSST Y1}
    \end{subfigure}
    \hspace{0.0\linewidth}
    \begin{subfigure}[b]{0.48\linewidth}
        \centering
        \includegraphics[width=\linewidth]{Figures/Y10/DENSITY.pdf}
        \caption{LSST Y10}
    \end{subfigure}
    \caption{Normalised PIT histograms for LSST Y1 (left panel) and Y10 (right panel) are depicted, under the baseline experiment. In each panel, the left subplots show lens samples, while the right subplots present source samples. The black solid lines represent the normalised PIT histogram of the whole sample, whereas the coloured solid lines illustrate true-redshift diagnostic intervals: five lens intervals of width \(0.3\) over \(0 < z_{\mathrm{true}} < 1.5\), and six source intervals of width \(0.5\) over \(0 < z_{\mathrm{true}} < 3.0\). Each PIT histogram uses ten quantile bins of width \(0.1\). The logarithmic axes display densities between \(0.08\) and \(8\); zero and out-of-range values are not shown. The upper row corresponds to the \texttt{Combination} dataset, used as an in-sample reference; the second row relates to the \texttt{Application} dataset trained on the augmented \texttt{Combination} dataset; and the bottom row illustrates the \texttt{Application} dataset trained on the \texttt{Degradation} dataset without data augmentation.}
    \label{fig:10}
\end{figure*}

\subsection{Limitations and future work}
One significant limitation of data augmentation for photometric‐redshift estimation is that the accuracy for galaxies at higher redshifts (\(z_\mathrm{true} \gtrsim 1.5\)) remains contingent upon the fidelity of the SED models in the mock catalogues even though our SOM-based weighted sample can largely align the distributions of galaxies in SED space. This shortcoming can be further alleviated by recalibrating the auto‐differentiable or semi‐analytic SED models against early LSST science observations, and by adopting a composite strategy that combines multiple simulations (e.g.\ \texttt{Cardinal}; \citealp{2024ApJ...961...59T} and \texttt{SKiLLS}; \citealp{2023A&A...670A.100L}), or even Stellar Population Synthesis (SPS) based mock galaxy SEDs using tools like \texttt{GalSBI}\citep{2025JCAP...06..007F, 2025A&A...703A.255T} or \texttt{pop-cosmos}\citep{2024ApJS..274...12A, 2024ApJ...975..145T, 2025ApJ...993..240T}, to achieve thorough coverage of the SOM space, thereby fully overlapping in the multi-dimensional colour-magnitude volume. By leveraging SOM‐based resampling, these comprehensive datasets can then be dynamically matched to observed photometry, offering a robust framework to mitigate residual discrepancies between simulations and real observations, particularly for faint, high‐redshift populations.  

Although mock catalogues such as \texttt{OpenUniverse2024} and \texttt{CosmoDC2} produce realistic galaxy properties, our analysis currently neglects pixel-level observational artefacts. Future work should incorporate forward-imaging simulations \citep{2023MNRAS.522.2801T}, for example employing \texttt{imSim} \citep{2014SPIE.9150E..14C} and \texttt{GalSim} \citep{2015A&C....10..121R}, to better capture systematics such as variable survey depth \citep{2016ApJS..226...24L, 2021A&A...648A..98J, 2025A&A...694A.259Y}, point-source contamination \citep{2015MNRAS.454.3500K}, and imperfect deblending of faint neighbouring sources \citep{2025arXiv250316680L}. These enhancements will enable a more realistic characterisation of photometric noise and selection biases, which is essential to ensure that data-augmentation techniques remain robust against the complexities inherent in real observational datasets.  

Our controlled‐degradation procedure provides only an approximation to the realistic spectroscopic samples anticipated for LSST cosmological analyses, since the distributions of magnitude and redshift limits will not be uniform across different surveys. Future work should therefore develop more comprehensive and physically motivated definitions of the spectroscopic selection function. Nonetheless, our approach furnishes a systematic framework for propagating selection-function effects into end-to-end photometric-redshift estimates. In practice, numerous spectroscopic surveys (e.g., DESI, PFS, and 4MOST), alongside deep narrow-band imaging programmes (e.g., COSMOS, PAUS, and J-PAS), will be integrated to optimise completeness and diversity within the multi-dimensional colour-magnitude space. Under this multi‐survey paradigm, controlled degradation remains a powerful tool: by selectively omitting or reweighting galaxy subsets from individual datasets, we can map the impact of incomplete spectroscopic coverage onto photometric redshifts, thereby enabling a rigorous characterisation of selection‐induced uncertainties.

Our PIT-based comparisons can be refined further in future work. Subdivision into true-redshift intervals can produce nonuniform PIT distributions even for an ideal conditional-density model, so we interpret interval-wise divergence losses as empirical indicators alongside the \texttt{Combination}-on-\texttt{Combination} reference and the point-estimate diagnostics. Complementary metrics that account explicitly for sample selection and true-redshift conditioning would help distinguish these effects from density-modelling errors.

In addition, as discussed in Section~\ref{sec:3}, our SOM-based data augmentation improves upon the simplified colour–magnitude–redshift boundaries used in \citet{2024ApJ...967L...6M} by providing a direct visualisation of the high-dimensional colour–magnitude manifold, thereby streamlining the identification of regions lacking spectroscopic coverage and optimising augmentation efficiency. Nonetheless, every dimensionality-reduction algorithm inevitably distorts the intrinsic topology. Future research will investigate cutting-edge techniques like Uniform Manifold Approximation and Projection \citep[UMAP;][]{2018arXiv180203426M} and \(t\)-distributed Stochastic Neighbour Embedding \citep[\(t\)-SNE;][]{vandermaatenJMLR:08a:v9}, aiming to improve augmentation fidelity. Moreover, more sophisticated and adaptable augmentation criteria can be defined within the joint colour-magnitude-redshift space, providing ample scope for further investigation.

A qualitative comparison with the approach utilised in \citet{2024ApJ...967L...6M} proves to be insightful. In that work, photometric redshifts were trained on a combination of currently available spectroscopic samples -- analogous to our LSST-Y1-like \texttt{Degradation} datasets -- but applied to an LSST-Y10 \texttt{Gold} sample, like our LSST Y10 \texttt{Application} datasets. Consequently, the level of missing spectroscopic information is considerably more significant than in our framework, where we anticipate broader spectroscopic coverage for Y10 analyses, though it remains anticipated. In \citet{2024ApJ...967L...6M}, the global photometric‐redshift metrics, such as biases, outlier fractions, and catastrophic rates, are reduced by a factor of \(\sim2\) when averaged over the full photometric sample. By contrast, in our analysis, these global statistics show only modest improvement, as they are dominated by low‐redshift galaxies with already high spectroscopic completeness. Nevertheless, our augmentation still yields clear benefits for the high‐redshift population. Taken together, the two configurations represent complementary scenarios: when the training sample is severely incomplete at the depth of the photometric data, augmentation improves overall performance; when the training sample is already broadly representative, the gains are concentrated at the faint, high‐redshift end.  

Given the differences between the mock catalogues and the varying photometric and spectroscopic selection functions adopted in this work and in \citet{2024ApJ...967L...6M}, a direct quantitative comparison is not straightforward. To enable a more rigorous assessment, we are actively analysing observational datasets from the \textit{Rubin} Data Preview~1 \citep[DP1\footnote{\url{https://dp1.lsst.io}},][]{10.71929/rubin/2570308}, which provides catalogue products over seven \(\sim 1 \,\mathrm{deg}^2\) fields for rigorous photometric-redshift validation. These analyses will benchmark our SOM-based augmentation against established methods -- both template-fitting and machine-learning -- to demonstrate its real-world efficacy. Further comparison using identical mock catalogues will enable careful evaluation of multiple augmentation strategies. The DP1 results will directly inform the integration of both colour–magnitude–redshift and SOM-based data augmentation workflows into the LSST DESC software ecosystem, with particular emphasis on the \texttt{RAIL} platform, thereby substantially enhancing photometric-redshift precision and robustness for Stage-IV cosmological studies.

Ultimately, this work has focused on constructing photometric and spectroscopic datasets and validating the efficacy of the SOM-based data-augmentation technique for photometric‐redshift estimation. In both Y1 and Y10 scenarios, the reduction of biases -- and in particular the significant decrease in catastrophic failure rates -- can substantially mitigate contamination when defining tomographic bins from photometric‐redshift estimates, thereby supporting the accurate calibration of ensemble redshift distributions. Building on this foundation, our forthcoming work will extend the framework to optimise tomographic binning strategies for both lens and source samples, following approaches discussed in the literature \cite{2021OJAp....4E..13Z, 2023ApJ...950...49M}, as well as to advance the calibration of ensemble redshift distributions. Leveraging the representative training of photometric redshifts and the improved performance achieved in both point estimates and conditional-density modelling, we will then enhance the calibration of ensemble redshift distributions by combining well‐behaved conditional PDFs with SOM‐based direct calibration \citep{2019MNRAS.489..820B, 2020A&A...637A.100W, 2025A&A...703A.144W}. This integrated approach aims to satisfy the stringent precision requirements of LSST cosmological analyses, particularly for weak gravitational lensing and large‐scale structure studies. Full details will be presented in the accompanying paper (Zhang et al., in preparation).  